\documentclass[11pt]{article}
\usepackage{amsmath,amsfonts}
\usepackage{graphicx}
%\usepackage{xcolor}
%\definecolor{gray}{rgb}{.75,.75,.75}
%\usepackage[landscape]{geometry}
%\usepackage{hyperref}
%\usepackage[bookmarks=true, pdfborder={0 0 .5}, urlbordercolor=gray]{hyperref}

\renewcommand{\thefootnote}{\fnsymbol{footnote}}


\addtolength{\topmargin}{-.7in}
\addtolength{\textheight}{1.3in}
\addtolength{\oddsidemargin}{-.4in}
\addtolength{\textwidth}{.8in}
\pagestyle{empty}

\newcommand{\ds}{\displaystyle}
\newcommand{\ts}{\textstyle}

\newcommand{\Z}{\mathbb{Z}}
\newcommand{\R}{\mathbb{R}}
\newcommand{\C}{\mathbb{C}}
\newcommand{\EE}{\mathcal{E}}
\newcommand{\tZ}{\tilde{\Z}}

\newcommand{\Sym}{\mathrm{Sym}}

\renewcommand{\circle}[1]{{\bigcirc\hspace{-.75em}#1}}

\begin{document}
\centerline{\bf Math 133\hspace{1.5em} \hfill{\Large\bf
Trigonometric Integrals}\hfill Stewart \S7.2}
\vspace{1em}

\noindent 
{\bf Products by substitution.} In this section, we develop methods
to find indefinite integrals (antiderivatives) of products of trig functions, beyond the
direct antiderivatives:
$$\ts
\int \cos(x)\,dx = \sin(x), \ \ 
\int \sin(x)\,dx = -\cos(x), \ \ 
\int \sec^2(x)\,dx = \tan(x), 
$$
\\[-2.4em]
$$\ts
\int \tan(x)\sec(x)\,dx = \sec(x),\ \ 
\int \csc^2(x)\,dx = -\cot(x), \ \ 
\int \cot(x)\csc(x)\,dx = -\csc(x).
$$ 
The easiest cases allow a simple trig substitution reducing to a polynomial, 
often using:
$$\sin^2(x)+\cos^2(x)=1, \qquad \tan^2(x)+1=\sec^2(x).$$
Let $m,n$ be integers with $n\geq 0$. The basic paradigms, with examples, are:

\begin{itemize}
\item[(a)] $\int  \sin^{m}(x) \cos^{2n+1}(x) \,dx$:
take $u=\sin(x)$, $du=\cos(x)\,dx$. \\
$\int\cos^{m}(x) \sin^{2n+1}(x)\,dx$:
take $u=\cos(x)$, $du=-\sin(x)\,dx$.  
$$\begin{array}{rcl}
\int\sin^{-10}(x) \cos^{3}(x)\,dx 
&=&
\int \sin^{-10}(x)\cos^{2}(x)\cdot \cos(x)\,dx
\\[.5em]&=&
\int \sin^{-10}(x)(1-\sin^2(x))\cdot \cos(x)\,dx
\\[.5em]
&=&
\int u^{-10}(1-u^2) \,du
\\[.5em]
&=&
\int u^{-10}-u^{-8} \,du
\\[.5em]
&=&
\tfrac{1}{-9}u^{-9} - \tfrac{1}{-7}u^{-7}+C
\\[.5em]
&=&
-\tfrac{1}{9}\sin^{-9}(x) + \tfrac{1}{7}\sin^{-7}(x)+C.
\end{array}$$
\\[-.9em]
$$\begin{array}{rcl}
\int \sin^{5}(x)\,dx 
&=&
-\int \sin^{4}(x)\cdot (-\sin(x))\,dx
\\[.5em]&=&
-\int (1-\cos^2(x))^2\cdot (-\sin(x))\,dx
\\[.5em]
&=&
-\int (1-u^2)^2 \,du
\\[.5em]
&=&
-\int (1-2u^2+u^4)\,du
\\[.5em]
&=&
-u + \tfrac23 u^3 - \tfrac1{5}u^{5}+C
\\[.5em]
&=&
-\cos(x) + \tfrac23\cos^{3}(x) - \tfrac1{5}\cos^{5}(x)+C.
\end{array}$$

\item[(b)] $\int \tan^{m}(x)\sec^{2n+2}(x)\,dx$: 
take $u=\tan(x)$, $du=\sec^2(x)\,dx$. 
$$\begin{array}{rcl}
\int \tan^{-3}(x)\sec^{6}(x)\,dx 
&=&
\int \tan^{-3}(x)\sec^{4}(x)\cdot \sec^2(x)\,dx 
\\[.5em]
&=&
\int \tan^{-3}(x)(\tan^2(x)+1)^2 \cdot \sec^2(x)\,du
\\[.5em]
&=&
\int u^{-3}(u^2+1)^2\,du
\\[.5em]
&=&
\int (u+2u^{-1}+u^{-3})\,du
\\[.5em]
&=&
\tfrac12 u^2 + \ln\!|u| - \tfrac12 u^{-2}+C
\\[.5em]
&=&
\tfrac12 \tan^2(x) +\ln|\!\tan(x)| -\tfrac12 \cot^2(x)+C.
\end{array}$$
 
\end{itemize}

\noindent
For brevity, we will henceforth omit the constant of integration $+C$.


\pagebreak

\noindent
{\bf Remaining cases.} What if a product of trig functions does not fit 
types (a) or (b)?

\begin{itemize}

\item  Rewrite in terms of $\sin$ and $\cos$, obtaining  
$\sin^{p}(x)\cos^q(x)$ for integers $p,q$.
Use type (a) if $p$ is odd and positive, or if $q$ is odd and positive.
Reduce to type (b) if $p$ and $q$ are both even with $p+q$ negative;
or if $p$ and $q$ are both odd and negative.\\
$$\hspace{-4.0em}
\int \sin^2(x)\cos(x)\tan^2(x)\csc(x)\,dx 
\ =\ \int \sin^2(x)\cos(x)\,\frac{\sin^2(x)}{\cos^2(x)}\,\frac{1}{\sin(x)}\,dx
$$
$$\hspace{3.0em}
\ =\ \int (\cos(x))^{-1}\sin^{3}(x)\,dx
\ =\  -\ln|\cos(x)|+\tfrac12\cos^2(x), \ \ 
\text{type (a)}.
$$
$$
\int \csc^2(x)\,dx 
\ =\ \int \frac{1}{\sin^2(x)}\,dx
\ =\ \int \frac{\cos^2(x)}{\sin^2(x)}\frac{1}{\cos^2(x)}\,dx
\ =\ \int \tan^{-2}(x)\sec^2(x)\,dx
$$
$$\hspace{3.0em}
\ =\ \int u^{-2}\,du
\ =\  -u^{-1} \ = \
\ = \ -(\tan(x))^{-1} \ = \ -\cot(x), \ \ 
\text{type (b)}.
$$
$$\hspace{-6.0em}
\int \frac{1}{\sin^3(x)\cos(x)}\,dx
\ =\ \int \frac{\cos^3(x)}{\sin^3(x)}\frac{1}{\cos^4(x)}\,dx
\ =\ \int \tan^{-3}(x)\sec^4(x)\,dx
$$
$$\hspace{3.0em}
\ =\ \int (u^2+1)u^{-3}\,du
\ =\  \ln|\!\tan(x)| -\tfrac12\cot^2(x), \ \ 
\text{type (b)}.
$$



\item Pythagorean identities: \quad
$\int \tan^4(x)\,dx\ =\ \int (\sec^2(x)-1)\tan^2(x)\,dx $\\
$\ =\ \int \sec^2(x)\tan^2(x)\,dx - \int (\sec^2(x)-1)\,dx \ =\ \frac13\tan^3(x)-\tan(x) + x$,
type (b).

\item For even positive powers of $\sin$ and $\cos$,
rewrite using the identities:\footnote{
Or: $\int\sin(x)\sin(x)\,dx 
\stackrel{\text{IBP}}= -\cos(x)\sin(x)-\int (-\cos(x))\cos(x)\,dx 
\stackrel{\text{Pyth}}= -\cos(x)\sin(x)+\int 1-\sin^2(x)\,dx$\\
$=-\cos(x)\sin(x)+x-\int\sin^2(x)\,dx$. Solve for $\int\sin^2(x)\,dx=\frac12(x-\cos(x)\sin(x))
=\frac12 x-\frac14\sin(2x)$.
}
\\[-.6em]
$$\sin^2(x)=\tfrac12(1{-}\cos(2x)), 
\quad
\cos^2(x)=\tfrac12(1{+}\cos(2x)),\quad
\sin(x)\cos(x)=\tfrac12\sin(2x).
$$
\\[-1.7em]
$$
\int \sin^6(x)\,dx 
\ =\ 
\int (\tfrac12(1{-}\cos(2x))^3\,dx
$$
$$
\ =\ 
\tfrac18\int 1- 3\cos(2x)+3\cos^2(2x)-\cos^3(2x)\ dx
$$
$$
\ =\ 
\tfrac18\int 1-3\cos(2x)+\tfrac32(1{+}\cos(4x))-\cos^3(2x)\ dx\,,
$$
where we used the binomial formula $(a{+}b)^3=a^3+3a^2b + 3ab^2+b^3$.
Now each term can be done directly or by type (a)
with the substitution $u=\sin(2x)$.

\end{itemize}
%\medskip

\noindent
{\bf Recalcitrant cases.} In case of odd negative powers of $\sin$ or $\cos$, 
we need special tricks.
\\[.5em]
{\sc example:} 
The integral of secant $\int \sec(x)\,dx$ was needed by map-makers in the 1600's, 
when Calculus was first developed. It calibrates stretching in the Mercator projection, 
which makes map directions match compass directions.
An amazing trick from \S6.2:
$$
\int\! \sec(x)\,dx
\ = \
\int\! \sec(x)\,
\frac{\sec(x)+\tan(x)}{\sec(x)+\tan(x)}
\,dx
$$
\\[-2em]
$$
\ = \
\int\! 
\frac{1}{\sec(x)+\tan(x)}\cdot\left(\sec^2(x)+\sec(x)\tan(x)\right)dx
\ =\
\int\! \frac{1}{u}\,du
$$
for $u=\sec(x)+\tan(x)$, 
$du=(\sec(x)\tan(x)+\sec^2(x))\,dx$.

\pagebreak

\noindent
Thus:
$$
\int \sec(x)\,dx \  =\ \ln\!|u|
\ =\ 
\ln|\!\sec(x){+}\tan(x)|.
$$
Another trick for this is to write $\int\! \sec(x)\,dx =\int \frac{1}{\cos^2(x)}\,\cos(x)\,dx$, and substitute $u=\sin(x)$ to get $\int\!\! \frac{1}{1-u^2}\,du$.
We will see how to integrate such rational functions in \S7.4.
\\[.5em]
{\sc example:} Here is a tricky integration by parts, in which we get back to the same integral we started with:
$$
\int\!\sec^3(x)\,dx
\ =\ 
\int\!\sec(x)(1+\tan^2(x))\,dx
\ =\ 
\int\! \sec(x)\,dx 
+ \int\! \underbrace{\tan(x)}_u\,\underbrace{\sec(x)\tan(x)\,dx}_{dv}
$$
$$
\ =\ \ln\!|\!\sec(x){+}\tan(x)| 
+ \underbrace{\tan(x)}_u\underbrace{\sec(x)}_v
- \int\! \underbrace{\sec(x)}_v\,\underbrace{\sec^2(x)\,dx}_{du}
$$
Since we have $\int\!\sec^3(x)\,dx$ on both sides
{\it with opposite signs}, we can solve for it to get:
$$
\int\!\sec^3(x)\,dx
\ =\ 
\tfrac12\!\left(\!\vphantom{\sum}
\ln\!\left|\vphantom{x^2}\!\sec(x){+}\tan(x)\right| 
\,+\,\tan(x)\sec(x)
\right).
$$
{\bf Trig integrals with inside coefficients.} Add together the angle addition identities:
$$
\cos(a{+}b)  \ = \ \cos(a)\cos(b) - \sin(a)\sin(b)
$$
$$
\cos(a{-}b) \ = \  \cos(a)\cos(b) + \sin(a)\sin(b)
$$
\\[-2em]
$$
\tfrac12(\cos(a{+}b) + \cos(a{-}b)) \ = \  \cos(a)\cos(b).
$$
This allows us to do integrals of the form:
$$
\int \cos(nx)\cos(mx)\,dx
\ =\
\int \tfrac12\cos(nx{+}mx) + \tfrac12\cos(nx{-}mx)\ dx
$$
$$
\ =\ \tfrac{1}{2(n{+}m)}\sin\!\left(\!\vphantom{A^2}(n{+}m)x\right)
+ \tfrac{1}{2(n{-}m)}\sin\!\left(\!\vphantom{A^2}(n{-}m)x\right)+C\,.
$$
\\[-.7em]
Similarly:
$$
\tfrac12(\cos(a{-}b) - \cos(a{+}b)) \ = \  \sin(a)\sin(b)
$$
$$
\tfrac12(\sin(a{+}b) + \sin(a{-}b)) \ = \  \sin(a)\cos(b).
$$
\\[0em]
{\bf Tangent half-angle substitution.} For a really messy trigonometric integral
like:
$$
\int \!\!\frac{\cos^2(x)\sin(x)-2\tan(x)+1}{\sec^3(x)+\sin^3(x)+3\cos(x)\sin(x)+5}\ dx,
$$
there is a general method in the extra notes,
the Geometric Trig Substitution:
$$
\sin(x) =  \frac{2t}{1+t^2}\,,
\quad \ \ 
\cos(x) = \frac{1-t^2}{1+t^2}\,, 
\quad \ \ 
%x = \arcsin\!\left(\!\tfrac{2t}{1+t^2}\!\right), 
%\quad
dx = \frac{2}{1+t^2}\,dt\,,
\quad \ \ 
t =\tan(\tfrac12 x) = \frac{\sin(x)}{1+\cos(x)}\,.
$$ 
This reduces the trig integral to the integral of a rational function, 
which can be done by Partial Fractions (\S7.4).



\end{document}

%%%%%%%%%%%  Picture  %%%%%%%%%%
\\[0em]
\centerline{
%\includegraphics[width=2in]
{1-8-Discont-iv.png}
\qquad \qquad 
%\includegraphics[width=2in]
{1-8-Discont-v.png} 
}
\vspace{0em}

\noindent
%%%%%%%%%%%%%%%%%%%%%%%%%

%%%%%%%%%%%  Table  %%%%%%%%%%
In fact, we get the following derivatives:
$$\begin{array}{|c||c|c|c|c|c|c|}
\hline &&&&&& \\[-.9em] 
f(x) & \sin(x) & \cos(x) & \tan(x)  & \sec(x) & \csc(x) & \cot(x)
\\[.2em] \hline &&&&&& \\[-.9em]
f'(x) &\cos(x) & -\sin(x) & \sec^2(x) & \tan(x)\sec(x) & -\cot(x)\csc(x) & -\csc^2(x)
\\[.1em] \hline
\end{array}$$
\\[1em]
%%%%%%%%%%%%%%%%%%%%%%%%%%



















